\documentclass[11pt,a4paper]{article}
\usepackage[utf8]{inputenc}
\usepackage[T1]{fontenc}
\usepackage[brazilian]{babel}
\usepackage{lmodern}
\usepackage{microtype}
\usepackage{geometry}
\geometry{top=2.2cm,bottom=2.2cm,left=2.4cm,right=2.4cm}
\usepackage{amsmath,amssymb}
\usepackage{booktabs}
\usepackage{graphicx}
\usepackage{caption}
\usepackage{subcaption}
\usepackage{siunitx}
\usepackage{hyperref}
\usepackage{xcolor}
\usepackage{listings}
\usepackage{float}
\usepackage{placeins}
\hypersetup{colorlinks=true,linkcolor=black,urlcolor=black,citecolor=black}
\sisetup{output-decimal-marker={,},group-separator={.},group-minimum-digits=4}
\setlength{\parindent}{1.2em}
\setlength{\parskip}{0.25em}
\captionsetup{font=small,labelfont=bf}
\lstset{basicstyle=\ttfamily\small,frame=single,breaklines=true,columns=fullflexible}

\title{
    \textbf{Processamento de Alto Desempenho (PAD)}\\[0.8em]
    \large Relatório de Trabalho 1 -- Distância entre Vetores\\[0.35em]
    \normalsize Implementação, otimização e análise experimental de desempenho
}

\author{Xavier Paulino Sebastião}
\date{Agosto de 2026}

\begin{document}
\maketitle

\begin{abstract}
Este trabalho avalia o cálculo da distância quadrática entre um vetor de referência e um conjunto de vetores reais e, em uma segunda etapa, compara a implementação original com uma versão otimizada. A versão de referência, denominada \texttt{baseline\_v3}, utiliza um laço escalar, enquanto a versão \texttt{optimized\_v4} explora SIMD AVX, desenrolamento do laço, quatro acumuladores vetoriais independentes e qualificadores \texttt{restrict}, mantendo execução \emph{single-thread} e as mesmas flags principais de compilação. A campanha comparativa considera oito tamanhos de vetor, de $T=32$ a $T=4096$, e mede as duas implementações sob as mesmas condições experimentais. A versão otimizada apresentou ganho em todos os tamanhos, com \emph{speedup} entre aproximadamente $1{,}32\times$ e $1{,}85\times$. O maior ganho ocorreu em $T=2048$, enquanto em $T=4096$ o \emph{speedup} diminuiu para cerca de $1{,}66\times$. Os resultados mostram que a otimização reduz de forma consistente o custo por elemento e aumenta a taxa de processamento, embora o ganho não permaneça constante em toda a faixa de tamanhos.
\end{abstract}

\section{Introdução}
O problema consiste em calcular, para cada vetor $x_i$ de um conjunto $X$ com $N$ vetores, a distância quadrática em relação a um vetor de referência $q$. Cada vetor contém $T$ valores reais, e a distância é definida por
\begin{equation}
D(q,x_i)=\sum_{j=0}^{T-1}(q_j-x_{ij})^2,
\qquad i=0,\ldots,N-1.
\end{equation}
Cada resultado é armazenado após o processamento do respectivo vetor. Mantendo $N$ constante e variando $T$, a quantidade de elementos processados por execução é $N\times T$ e o trabalho aritmético cresce linearmente com $T$.

A primeira etapa do trabalho caracterizou o comportamento da implementação escalar para diferentes valores de $T$. A segunda etapa introduz uma versão otimizada do mesmo cálculo e avalia quanto essa otimização altera o tempo de execução, as taxas de processamento, o custo por elemento e a variabilidade. Em avaliações de desempenho, múltiplas observações e a caracterização da dispersão são importantes para evitar conclusões baseadas em execuções isoladas~\cite{hoefler2015,curtsinger2013,kalibera2013rigorous}.

\section{Implementações avaliadas}
\subsection{Versão base v3}
A implementação de referência foi desenvolvida em C++~\cite{schildt1998cpp}, com dados em precisão dupla e armazenamento contíguo da matriz $X$. Para cada vetor, o kernel percorre sequencialmente os $T$ elementos e acumula a soma dos quadrados das diferenças:

\begin{lstlisting}[language=C++]
for (size_t i = 0; i < N; ++i) {
    const double* row = X + i * T;
    double sum = 0.0;
    for (size_t j = 0; j < T; ++j) {
        const double d = q[j] - row[j];
        sum += d * d;
    }
    distances[i] = sum;
}
\end{lstlisting}

Essa versão é mantida no código da Parte 2 sem alteração funcional e é identificada como \texttt{baseline\_v3}. A manutenção do kernel original dentro do mesmo executável permite comparar as duas implementações na mesma campanha, evitando utilizar apenas medições obtidas em momentos distintos.

\subsection{Versão otimizada v4}
A versão otimizada preserva o mesmo layout de dados, o mesmo cálculo matemático e o mesmo orçamento de recursos: a execução continua \emph{single-thread}. A otimização foi realizada no laço interno e combina quatro técnicas. Primeiro, os ponteiros do kernel são declarados com \texttt{\_\_restrict\_\_}, reduzindo a ambiguidade de aliasing para o compilador. Segundo, quando AVX está disponível, quatro valores \texttt{double} são processados por instrução vetorial. Terceiro, o laço é desenrolado em blocos de 16 elementos. Quarto, são usados quatro acumuladores vetoriais independentes, diminuindo a cadeia serial de dependência da redução.

O trecho principal da versão AVX é representado por:

\begin{lstlisting}[language=C++]
__m256d acc0 = _mm256_setzero_pd();
__m256d acc1 = _mm256_setzero_pd();
__m256d acc2 = _mm256_setzero_pd();
__m256d acc3 = _mm256_setzero_pd();

for (; j + 15 < T; j += 16) {
    // quatro blocos de 4 doubles
    // load q, load x, subtracao, multiplicacao e acumulacao
}
\end{lstlisting}

Após o bloco principal, vetores AVX completos remanescentes são processados em grupos de quatro elementos e um laço escalar trata o restante. Assim, a implementação não depende de $T$ ser múltiplo de 16. A ordem da soma difere da versão escalar; por isso, a validação de corretude compara as duas implementações com tolerância numérica, em vez de exigir igualdade bit a bit. A compilação mantém C++20, \texttt{-O3} e \texttt{-march=native}, sem \texttt{-ffast-math}. Otimizações SIMD e transformações de laços são mecanismos diretamente relacionados ao código gerado e ao uso da microarquitetura~\cite{intelOpt2026}.

\section{Metodologia experimental}
\subsection{Plataforma e ambiente}
As aferições foram executadas em um Lenovo ThinkPad T530 equipado com processador Intel Core i5-3320M a \SI{2.60}{\giga\hertz}, arquitetura x86-64, família 6, modelo 58. A CPU disponibiliza AVX, SSE4.1 e SSE4.2. A compilação foi realizada com GCC 10.5.0, usando \texttt{/usr/bin/g++}. O sistema operacional utilizado foi GNU/Linux Ubuntu 24.04.4 LTS. Os scripts de análise e geração de gráficos foram executados no ambiente Conda \texttt{t0\_vector\_distance}.

\subsection{Campanha comparativa}
A comparação da Parte 2 mede \texttt{baseline\_v3} e \texttt{optimized\_v4} novamente na mesma campanha. Para cada bloco experimental, as 16 combinações formadas por duas implementações e oito valores de $T$ são ordenadas de forma pseudoaleatória e determinística. Essa estratégia reduz a associação sistemática entre uma variante e o momento temporal da execução, aspecto relevante porque condições aparentemente secundárias da configuração experimental podem introduzir viés em medições de desempenho~\cite{mytkowicz2009}.

Na campanha representada pelos resultados comparativos, $N=8192$ permaneceu constante para as duas implementações e para todos os valores de
\[
T\in\{32,64,128,256,512,1024,2048,4096\}.
\]
Cada combinação $(\text{implementação},T)$ foi aferida repetidamente, e as estatísticas apresentadas utilizam a mediana como medida central. O uso de repetições permite caracterizar a variabilidade e reduzir a dependência de uma única observação~\cite{kalibera2013rigorous}.

A região cronometrada contém exclusivamente o cálculo e o armazenamento das $N$ distâncias. Alocação, geração e inicialização dos dados, aquecimento, checksum, escrita de arquivos e apresentação dos resultados permanecem fora do intervalo medido. A temporização utiliza \texttt{CLOCK\_MONOTONIC\_RAW}~\cite{buttazzo2013ptask}.

\subsection{Índices de desempenho}
Além do tempo de execução $t$, foram calculadas as taxas
\begin{equation}
P_v=\frac{N}{t},
\qquad
P_e=\frac{N\,T}{t},
\end{equation}
onde $P_v$ representa vetores processados por segundo e $P_e$ elementos processados por segundo. O custo temporal por elemento é dado por
\begin{equation}
C_e=\frac{t\times 10^9}{N\,T}\quad\text{ns/elemento}.
\end{equation}
Para quantificar diretamente o efeito da otimização, foi calculado
\begin{equation}
S(T)=\frac{t_{\mathrm{baseline}}(T)}{t_{\mathrm{optimized}}(T)},
\end{equation}
em que $S(T)>1$ representa ganho de desempenho. A redução percentual do tempo é
\begin{equation}
R(T)=100\left(1-\frac{t_{\mathrm{optimized}}(T)}{t_{\mathrm{baseline}}(T)}\right).
\end{equation}

\section{Resultados}
\subsection{Resumo quantitativo}
A Tabela~\ref{tab:comparacao} reúne as principais métricas das duas implementações. A versão otimizada foi mais rápida em todos os tamanhos avaliados. O ganho não é uniforme: o \emph{speedup} cresce de aproximadamente $1{,}32\times$ em $T=32$ até $1{,}85\times$ em $T=2048$ e diminui para cerca de $1{,}66\times$ em $T=4096$.

\begin{table}[H]
\centering
\caption{Comparação entre a implementação baseline v3 e a optimized v4.}
\label{tab:comparacao}
\scriptsize
\resizebox{\textwidth}{!}{%
\begin{tabular}{rrrrrrrrr}
\toprule
$T$ & \multicolumn{2}{c}{Tempo mediano (ms)} & \multicolumn{2}{c}{$P_e$ ($10^9$ elem./s)} & \multicolumn{2}{c}{ns/elem.} & Speedup & Redução (\%) \\
\cmidrule(lr){2-3}\cmidrule(lr){4-5}\cmidrule(lr){6-7}
 & v3 & v4 & v3 & v4 & v3 & v4 &  &  \\
\midrule
32   & 0,198 & 0,150 & 1,325 & 1,752 & 0,755 & 0,571 & 1,32 & 24,4 \\
64   & 0,463 & 0,290 & 1,132 & 1,810 & 0,883 & 0,553 & 1,60 & 37,4 \\
128  & 1,015 & 0,602 & 1,033 & 1,742 & 0,968 & 0,574 & 1,69 & 40,7 \\
256  & 2,111 & 1,190 & 0,993 & 1,762 & 1,007 & 0,568 & 1,77 & 43,6 \\
512  & 4,225 & 2,333 & 0,993 & 1,798 & 1,007 & 0,556 & 1,81 & 44,8 \\
1024 & 8,570 & 4,742 & 0,979 & 1,769 & 1,022 & 0,565 & 1,81 & 44,7 \\
2048 & 17,247 & 9,305 & 0,973 & 1,803 & 1,028 & 0,555 & 1,85 & 46,1 \\
4096 & 34,783 & 20,917 & 0,965 & 1,604 & 1,037 & 0,623 & 1,66 & 39,9 \\
\bottomrule
\end{tabular}%
}
\end{table}

\subsection{Tempo de execução}
A Figura~\ref{fig:tempo-comp} mostra que o tempo cresce com $T$ nas duas versões, mas a curva otimizada permanece abaixo da baseline em toda a faixa. Em $T=32$, o tempo mediano é aproximadamente \SI{0.198}{\milli\second} na v3 e \SI{0.150}{\milli\second} na v4. Em $T=4096$, os valores são aproximadamente \SI{34.783}{\milli\second} e \SI{20.917}{\milli\second}, respectivamente.

\begin{figure}[H]
\centering
\includegraphics[width=0.82\linewidth]{execution_time_comparison.png}
\caption{Tempo mediano de execução das versões baseline v3 e optimized v4.}
\label{fig:tempo-comp}
\end{figure}

\subsection{Taxas de processamento}
A taxa de vetores por segundo aumenta com a otimização em todos os tamanhos, como mostra a Figura~\ref{fig:vet-comp}. A diferença é especialmente visível nos tamanhos intermediários, nos quais a v4 conclui aproximadamente entre $1{,}7\times$ e $1{,}9\times$ mais vetores por unidade de tempo.

\begin{figure}[H]
\centering
\includegraphics[width=0.80\linewidth]{vectors_per_second_comparison.png}
\caption{Taxa mediana de processamento de vetores.}
\label{fig:vet-comp}
\end{figure}

A Figura~\ref{fig:elem-comp} apresenta a taxa de elementos, que é mais adequada para comparar o custo por unidade de trabalho. A baseline varia de aproximadamente 1,33 bilhão de elementos/s em $T=32$ para 0,97 bilhão em $T=4096$. A versão otimizada permanece próxima de 1,75--1,81 bilhão de elementos/s entre $T=32$ e $T=2048$, caindo para aproximadamente 1,60 bilhão de elementos/s em $T=4096$.

\begin{figure}[H]
\centering
\includegraphics[width=0.80\linewidth]{elements_per_second_comparison.png}
\caption{Taxa mediana de processamento de elementos.}
\label{fig:elem-comp}
\end{figure}

\subsection{Custo por elemento}
O custo por elemento evidencia diretamente o efeito da otimização. Na baseline, esse custo cresce de aproximadamente 0,755 ns/elemento em $T=32$ para 1,037 ns/elemento em $T=4096$. Na v4, permanece entre aproximadamente 0,553 e 0,574 ns/elemento de $T=32$ a $T=2048$, com aumento para cerca de 0,623 ns/elemento no maior tamanho.

\begin{figure}[H]
\centering
\includegraphics[width=0.80\linewidth]{ns_per_element_comparison.png}
\caption{Custo mediano por elemento nas duas implementações.}
\label{fig:nselem-comp}
\end{figure}

A estabilidade relativa do custo por elemento na v4 entre $T=64$ e $T=2048$ indica que a nova organização do laço reduz parte do custo observado na implementação escalar. Entretanto, a elevação em $T=4096$ mostra que o ganho não é independente do tamanho da entrada. Como o experimento não coleta contadores de hardware ou largura de banda de memória, não é possível identificar de forma conclusiva se essa redução do ganho decorre de cache, memória principal, prefetch, frequência da CPU ou outro fator microarquitetural. A infraestrutura \texttt{perf\_event} do Linux permitiria medir eventos como ciclos, instruções e faltas de cache, mas esses indicadores não fazem parte dos dados apresentados~\cite{linuxPerf2026}.

\subsection{Speedup}
A Figura~\ref{fig:speedup} resume o efeito da otimização. O ganho cresce de $1{,}32\times$ em $T=32$ para aproximadamente $1{,}60\times$ em $T=64$ e supera $1{,}77\times$ a partir de $T=256$. O maior valor observado é aproximadamente $1{,}85\times$ em $T=2048$. Em $T=4096$, o \emph{speedup} cai para aproximadamente $1{,}66\times$, mas permanece substancialmente acima de 1.

\begin{figure}[H]
\centering
\includegraphics[width=0.80\linewidth]{speedup.png}
\caption{Speedup da versão optimized v4 em relação à baseline v3. A linha tracejada representa ausência de ganho.}
\label{fig:speedup}
\end{figure}

\subsection{Variabilidade}
A Figura~\ref{fig:var-comp} apresenta as distribuições das aferições. As duas versões mostram dispersão relativamente pequena na maior parte dos tamanhos, com observações superiores isoladas em alguns casos. Em $T=4096$, a baseline apresenta outliers acima do núcleo principal da distribuição; a versão otimizada também apresenta observações elevadas, porém em uma escala de tempo menor. As observações foram preservadas, evitando reduzir artificialmente a variabilidade da campanha.

\begin{figure}[H]
\centering
\includegraphics[width=0.82\linewidth]{execution_time_variability_comparison.png}
\caption{Distribuição dos tempos de execução das duas implementações.}
\label{fig:var-comp}
\end{figure}

\section{Discussão}
Os resultados mostram que a otimização produz ganho consistente em toda a faixa de tamanhos, mas não um fator constante. Esse comportamento é compatível com a natureza das mudanças realizadas no kernel. A versão escalar mantém uma única variável de acumulação, de modo que cada atualização de \texttt{sum} depende do valor produzido na iteração anterior. A v4 utiliza quatro acumuladores AVX independentes no bloco principal, permitindo maior paralelismo em nível de instrução e processando múltiplos elementos simultaneamente. O desenrolamento também reduz a frequência relativa do controle do laço.

O ganho de aproximadamente $1{,}8\times$ nos tamanhos intermediários e grandes é inferior ao fator quatro correspondente à largura de um registrador AVX com valores \texttt{double}. Isso não representa uma anomalia. O tempo total do kernel não é determinado apenas pela quantidade de operações aritméticas: carregamentos, armazenamento do resultado, redução horizontal, controle de laço, latências de instruções e movimento de dados também participam do custo. Além disso, a implementação já era compilada com \texttt{-O3 -march=native}, de modo que parte das otimizações de baixo nível podia ser explorada pelo compilador mesmo antes da introdução explícita do kernel AVX.

A queda de \emph{speedup} em $T=4096$ merece atenção. Nesse ponto, tanto o throughput da v4 quanto o ganho relativo diminuem, enquanto o custo por elemento aumenta de aproximadamente 0,555 ns em $T=2048$ para 0,623 ns em $T=4096$. O resultado demonstra uma mudança mensurável, mas não identifica sua causa. A hierarquia de memória e o prefetch de hardware podem influenciar percursos sequenciais desse tipo~\cite{intelOpt2026}; contudo, sem contadores de cache, tráfego de memória ou frequência instantânea, associar a queda de ganho especificamente a um desses mecanismos iria além da evidência coletada.

A comparação também foi estruturada para reduzir um problema metodológico presente em uma comparação simples entre duas campanhas independentes. Baseline e versão otimizada foram medidas novamente no mesmo período experimental e intercaladas em blocos, reduzindo a possibilidade de que uma diferença temporal da máquina seja confundida com o efeito da otimização. Essa prática é coerente com recomendações de benchmarking que enfatizam repetição, controle de variáveis e análise explícita de variabilidade~\cite{hoefler2015,kalibera2013rigorous,curtsinger2013}.

\section{Conclusão}
A otimização da implementação produziu melhoria de desempenho em todos os tamanhos avaliados sem recorrer a paralelismo entre cores. A versão optimized v4 manteve o mesmo cálculo, o mesmo layout de dados e execução \emph{single-thread}, concentrando as mudanças no kernel por meio de AVX, desenrolamento do laço, acumuladores independentes e \texttt{restrict}.

O \emph{speedup} variou de aproximadamente $1{,}32\times$ a $1{,}85\times$. O maior ganho foi observado em $T=2048$, com redução de tempo de aproximadamente 46,1\%. Em $T=4096$, o ganho caiu para cerca de $1{,}66\times$, ainda correspondendo a redução de aproximadamente 39,9\% no tempo mediano. A taxa de elementos da versão otimizada permaneceu próxima de 1,8 bilhão de elementos/s na maior parte da faixa e o custo por elemento ficou próximo de 0,56 ns até $T=2048$.

Os resultados confirmam que a reorganização do kernel reduz de forma substancial o custo computacional da operação, mas também mostram que a eficiência da otimização depende do tamanho da entrada. A queda relativa observada no maior caso não pode ser explicada de forma conclusiva apenas pelos tempos coletados. Assim, a evidência experimental sustenta o ganho da versão otimizada, mas não uma atribuição causal específica para a mudança de comportamento em $T=4096$.

\begin{thebibliography}{99}

\bibitem{buttazzo2013ptask}
G. Buttazzo and G. Lipari,
``Ptask: An educational C library for programming real-time systems on Linux,''
in \emph{2013 IEEE 18th Conference on Emerging Technologies \& Factory Automation (ETFA)},
pp. 1--8, 2013.

\bibitem{hoefler2015}
T. Hoefler and R. Belli,
``Scientific benchmarking of parallel computing systems: Twelve ways to tell the masses when reporting performance results,''
in \emph{Proceedings of the International Conference for High Performance Computing, Networking, Storage and Analysis},
pp. 1--12, 2015.

\bibitem{curtsinger2013}
C. Curtsinger and E. D. Berger,
``Stabilizer: Statistically sound performance evaluation,''
\emph{ACM SIGARCH Computer Architecture News},
vol. 41, no. 1, pp. 219--228, 2013.

\bibitem{intelOpt2026}
Intel Corporation,
\emph{Intel 64 and IA-32 Architectures Optimization Reference Manual}, Version 050, 2026.
Disponível em: \url{https://www.intel.com/content/www/us/en/developer/articles/technical/intel64-and-ia32-architectures-optimization.html}.

\bibitem{linuxPerf2026}
M. Kerrisk and Linux man-pages contributors,
``perf\_event\_open(2) -- set up performance monitoring,'' \emph{Linux man-pages}, 2026.
Disponível em: \url{https://man7.org/linux/man-pages/man2/perf_event_open.2.html}.

\bibitem{mytkowicz2009}
T. Mytkowicz, A. Diwan, M. Hauswirth, and P. F. Sweeney,
``Producing wrong data without doing anything obviously wrong!,''
\emph{ACM SIGPLAN Notices},
vol. 44, no. 3, pp. 265--276, 2009.

\bibitem{schildt1998cpp}
H. Schildt,
\emph{C++: The Complete Reference}.
McGraw-Hill, 1998.

\bibitem{kalibera2013rigorous}
T. Kalibera and R. Jones,
``Rigorous benchmarking in reasonable time,''
in \emph{Proceedings of the 2013 International Symposium on Memory Management (ISMM)},
pp. 63--74, 2013.

\end{thebibliography}

\end{document}
